% Fuente editable reutilizable. El DOCX en recursos/ conserva el original recibido.

\begin{enumerate}[leftmargin=*, itemsep=12pt]

\item En sus propias palabras, defina que es la fonética en un sentido amplio, es decir, dando una definición que incluya a las lenguas de señas. En su definición, no olvide detallar las subramas de la fonética e ilustre comparativamente cada rama para las lenguas orales y las lenguas de señas.

\item Explore los siguientes materiales: \href{https://www.youtube.com/playlist?list=PLh1FSFqMn5hATtH\_meS1I8qM8H9pJwpr2}{Recontar Historias: Rana, ¿Dónde Estás?} y/o \href{https://www.youtube.com/playlist?list=PLh1FSFqMn5hDMJ5356GeYGkuy7-ucvV8j}{Recontar Historias: La Historia del Niño}. Para cada material seleccione una seña y explore tres diferentes realizaciones fonéticas de la misma. Luego, haga una revisión de si estas producciones son fonéticamente idénticas (o no) en las diferentes ocurrencias, presentando las respectivas descripciones articulatorio-perceptuales. Tome por ejemplo, el caso para la seña \href{https://drive.google.com/file/d/1B1Z5tOB7Jnn7iWsXKE9Xr1jVaF9wBkFd/view?usp=drive\_link}{CASA} (\href{https://www.youtube.com/watch?t=37\&v=Lhv2ZGsdzrM}{0.37}, \href{https://www.youtube.com/watch?t=41\&v=izfaNu3DkNg}{0.41}, \href{https://www.youtube.com/watch?t=16\&v=vWmA1u4-N5o}{0.16}), cada una articulada con diferencias fonéticas entre sí en los materiales de \href{https://www.youtube.com/playlist?list=PLh1FSFqMn5hDMJ5356GeYGkuy7-ucvV8j}{Recontar Historias: La Historia del Niño}.


NOTA.



Al desarrollar este punto no olvide incluir en el desarrollo del taller imágenes de la seña en el discurso seleccionado, así como el enlace al vídeo y la marca de tiempo en la que ocurre. Asimismo, no olvide que la descripción fonética se hace en términos físicos y anatómicos.


\item En sus propias palabras, defina que es la fonología en un sentido amplio, es decir, dando una definición que incluya a las lenguas de señas. En su definición, no olvide detallar las unidades de análisis básicas de la fonología presentadas en clase e ilustre estas en pares mínimos. Asimismo, señale de manera explícita la relación —y la diferencia— entre la fonología y la fonética.

\item Examine los siguientes pares de señas. Luego, describa qué cambio fonológico identifica para cada caso (haga caso omiso a los RNMs).

\begin{center}\begin{tabular}{ll}

\href{https://drive.google.com/file/d/1AVrkiZDxv0dDYBSvmPsWHvxJtMWPr0lO/view?usp=drive\_link}{TORO-A} & \href{https://drive.google.com/file/d/1AVsL2i6q\_lynJ7EGMli4vbLgcGlXLsBk/view?usp=drive\_link}{TORO-B} \\[5pt]

\href{https://drive.google.com/file/d/1At4c1E5j9YGr468nGAW5eG5hCORmcsYD/view?usp=drive\_link}{TREN-A} & \href{https://drive.google.com/file/d/1AkfI6k6QHorgac2BuR\_znuCb5VfMtP77/view?usp=drive\_link}{TREN-B} \\[5pt]

\href{https://drive.google.com/file/d/1AXutZpCEkWHc62\_JArNaTR32TfsoUw5c/view?usp=drive\_link}{MAZORCA-A} & \href{https://drive.google.com/file/d/1A\_Ct3E7so1UZRenrKTgDvzpTRbm88kJG/view?usp=drive\_link}{MAZORCA-B} \\[5pt]

\href{https://drive.google.com/file/d/1ATcgn3aEv0bko\_VM-DGry2yRKLve4ygj/view?usp=drive\_link}{MARIPOSA-A} & \href{https://drive.google.com/file/d/1AMlQWaGYu3kwJStCfA91HZHNdrxzcZyw/view?usp=drive\_link}{MARIPOSA-B} \\[5pt]

\end{tabular}\end{center}

\item Elabore una reflexión de no menos de una cuartilla y no más de una y media, donde se desarrolle la importancia de los temas vistos en clase en dos frentes: el de la importancia de estos temas para el entendimiento y posicionamiento de las lenguas de señas, entre estas la LSC; y ii) el de la importancia de estos temas para el proceso formativo como futuros intérpretes-traductores de LSC-Español.

\end{enumerate}
